% Method (one paragraph per line).
\section{Detector, samples and method}
\label{sec:method}

\subsection{Detector and samples}
\label{sec:detector}

The detector has 22 silicon layers arranged in a barrel and two endcaps, inside a solenoidal magnetic field. The innermost layers are pixel sensors, and the outer layers are strip sensors with a coarser resolution along the beam axis. Figure~\ref{fig:detector} shows the layout in the transverse and longitudinal views. Each charged particle leaves one hit per layer it crosses, plus occasional extra hits where modules overlap. Hits are recorded as three-dimensional space points in cylindrical coordinates.

\begin{figure}[tbp]
  \centering
  \fbox{\parbox[c][0.45\linewidth][c]{0.8\linewidth}{\centering placeholder artwork: detector}}
  \caption{Layout of the simulated detector in the transverse (left) and longitudinal (right) views. Pixel layers are drawn in dark grey and strip layers in light grey.}
  \label{fig:detector}
\end{figure}

We use three samples of simulated events. The generic sample contains 10,000 top-quark pair events without pile-up. The signal sample contains 10,000 events in which a heavy neutral lepton decays to a muon and two other charged particles\footnote{The decay length of each heavy neutral lepton is drawn from an exponential distribution with the nominal lifetime.}. The external test sample was generated with a different lepton mass and lifetime after the design of the study was fixed. Each sample is split into 8,000 training, 1,000 validation and 1,000 test events; Table~\ref{tab:samples} lists them. Selection of anything (models, thresholds, working points, bins) uses validation events only.

Events are generated with a standard parton-shower generator and passed through a fast simulation of the detector. The simulation propagates each charged particle through the magnetic field and the material of every layer. Multiple scattering and energy loss are applied layer by layer, and hits are smeared with the resolution of the sensor type. Noise hits are not simulated, because their rate in silicon detectors is small compared with the hits from particles. Secondary particles from interactions with the material are kept and treated like any other particle. They are a realistic source of displaced tracks in the generic sample.

The heavy neutral lepton is produced in the decay of a W boson together with a prompt muon. It travels a distance set by its lifetime and then decays to a second muon and a pair of quarks or leptons. The charged decay products therefore start at the position of the decay, which can be anywhere inside the tracker. We use the prompt muon of the same event as a reference that both trackers should find. Its efficiency is close to one for both trackers, which shows that the comparison is not affected by a gross failure of either.

Before the graph is built, the hits of each event are sorted by layer and by azimuthal angle. Hits from the same module that lie closer than the sensor pitch are merged into one space point. The merged space point keeps the identity of the particle that left most of its charge. Space points are written to disk together with the truth information needed for training and evaluation. The truth information never reaches the network at inference time.

Events are processed one at a time, without batching across events at inference. Each event is independent, so the pipeline can be distributed over many processors without communication. The same code runs the training, the validation and the evaluation, with only the input files changed. This guarantees that the two trackers see the same preprocessing.

\begin{table}[tbp]
  \caption{Simulated samples and their splits. The external test sample is used only for the final comparison.}
  \label{tab:samples}
  \centering
  \begin{tabular}{lrrr}
    \hline\hline
    Sample & Training & Validation & Test \\
    \hline
    Top-quark pairs & 8,000 & 1,000 & 1,000 \\
    Heavy neutral lepton & 8,000 & 1,000 & 1,000 \\
    External signal & --- & --- & 1,000 \\
    \hline\hline
  \end{tabular}
\end{table}

A target particle has $p_T > 1$~GeV, $|\eta| < 4$ and at least five space points. A displaced particle is a target with a transverse impact parameter $|d_0| > 2$~mm. The impact parameter is defined from the particle's own hits by a circle fit in the transverse plane,
\begin{equation}
  d_0 = \left(\vec{p} \times \hat{t}\right)_z ,
  \label{eq:d0}
\end{equation}
where $\vec{p}$ is the point of closest approach to the beam axis and $\hat{t}$ the unit direction of motion at that point.

\subsection{Graph and network}
\label{sec:graph}

The pipeline has four stages: graph construction, edge classification, edge selection and track building. Figure~\ref{fig:pipeline} shows the stages and the information that flows between them. Graph construction connects each hit to the hits on adjacent layers that a particle could plausibly reach. A module map, learned from the training events, lists the pairs of modules that particles actually connect\footnote{The module map is built from training events only, so it contains no information about the test sample.}. Only hit pairs with $|\Delta z| < 300$~mm on adjacent layers enter the graph. On average a graph holds 43,260 nodes and 146,892 edges per event. It keeps 98.23\% of the true edges at an edge purity of 0.42.

The edge classifier is an interaction network that updates node and edge features in alternating steps. Each of the 8 message-passing steps lets information travel one edge further along a candidate track. The node and edge networks are multilayer perceptrons with 128 hidden units each. The full model has 1,344,392 trainable parameters. After the last step, a final perceptron assigns each edge a score between zero and one.

\begin{figure*}[tbp]
  \centering
  \fbox{\parbox[c][0.45\linewidth][c]{0.9\textwidth}{\centering placeholder artwork: pipeline}}
  \caption{The four stages of the pipeline. Graph construction and track building are the same for both trackers; only the training events of the edge classifier differ.}
  \label{fig:pipeline}
\end{figure*}

The network is trained with a weighted binary cross-entropy,
\begin{equation}
  \mathcal{L} = -\sum_{e} w_e \left[ y_e \log s_e + (1 - y_e) \log (1 - s_e) \right] ,
  \label{eq:loss}
\end{equation}
where $s_e$ is the score of edge $e$, $y_e$ is one for a true edge and zero otherwise, and the weight $w_e$ balances the rare true edges against the many false ones.

\subsection{Training, track building and metrics}
\label{sec:metrics}

We train for 60 epochs with the Adam optimizer~\cite{Bauer2016Adam}, a batch of 16 events and a learning rate of $3\times10^{-4}$. Every configuration is trained with 3 random seeds, and the reported numbers are averages over the seeds. The two trackers differ only in their training events: generic events for one, signal events for the other. Training stops at the epoch with the best validation efficiency, judged on validation events of the training process itself. No test event influences any choice of the pipeline.

Edges with a score above one half are kept. The threshold was fixed on validation events of the generic sample and is the same for both trackers. Track building then follows the kept edges from each hit with no incoming edge to the hit with no outgoing edge. Where a path branches, the builder follows the edge with the higher score and starts a new candidate at the other. A candidate becomes a track when it holds at least five hits.

A reconstructed track matches a particle if more than half of its hits come from that particle and more than half of the particle's hits are on the track. This criterion is called double majority. The efficiency is the fraction of target particles matched by a track,
\begin{equation}
  \epsilon = \frac{N_{\text{matched}}}{N_{\text{targets}}} ,
  \label{eq:eff}
\end{equation}
where both counts use the same population of targets; Eq.~(\ref{eq:eff}) is evaluated separately for each tracker.

A fake track is a track with at least five hits that matches no particle\footnote{A track may match at most one particle and a particle at most one track under double majority.}. The fake-track rate is the fraction of such tracks among all tracks with at least five hits. A duplicate is a second track matched to a particle that already has one. The three numbers together describe what a search would see: the particles it finds, the noise it must reject and the double counting it must undo.

The headline number is the efficiency difference between the trackers on displaced targets,
\begin{align}
  \Delta &= \epsilon_{\text{generic}} - \epsilon_{\text{signal}} , \label{eq:delta} \\
  \sigma_\Delta^2 &= \sigma_{\text{boot}}^2 + \sigma_{\text{seeds}}^2 . \label{eq:sigma}
\end{align}
with the two components defined in Appendix~\ref{app:stat}.

Its uncertainty combines two sources in quadrature: the bootstrap of the test events and the spread over training seeds~\cite{Novak2023Bootstrap}. The bootstrap draws the same resampled event lists for both trackers, so their difference is not affected by which events happen to be in the test set~\cite{Quinn2016Statistics}. Appendix~\ref{app:stat} gives the details of both components. Both trackers are evaluated on exactly the same test events and the same target particles.

\subsection{Track fit and validation}
\label{sec:validation}

After track building, every track is fitted with a helix to obtain its parameters. The fit uses only the hits on the track and does not depend on which tracker built it. A track whose fit fails is kept for the efficiency but marked in the resolution study. Fits fail mostly for tracks that combine hits of two particles, which are also the tracks that fail double majority. The fitted parameters enter only the resolution study in the results.

We validated every stage of the pipeline on validation events before looking at the test events. The graph construction was checked by counting the true edges it keeps, layer pair by layer pair. The edge classifier was checked by comparing its score distribution for true and false edges between training and validation events. The track builder was checked on graphs that contain only true edges, where it must reproduce the true particles. All checks passed with the same settings for both trackers. We then froze the pipeline, recorded its configuration and only afterwards evaluated the test events.

Training and inference run on one graphics processor per job. A full training run of either tracker takes less than a day. The evaluation of all test samples for one network takes a few hours. All runs use the same software versions, which are recorded with their outputs. Floating-point results differ slightly between processors, but by far less than the seed spread.
